\ifx\danslelivre\undefined
\documentclass[10pt]{../fichiers-configuration-latex/smfart}
\usepackage{ifthen}
\newboolean{english}
\newboolean{francais}
\setboolean{english}{true}
\setboolean{francais}{false}
\input{../fichiers-configuration-latex/paquets}
\selectlanguage{english}
\input{../fichiers-configuration-latex/lemme-proposition-theoreme}
\input{../fichiers-configuration-latex/numerotation}
\input{../fichiers-configuration-latex/macros}

\externaldocument{../00-introduction/introduction}
\externaldocument{../01-quasi-excellents/excellent}
\externaldocument{../02-topologies/topologies}
\externaldocument{../03-approximation/approximation}
\externaldocument{../04-Cohen-Gabber/Cohen-Gabber}
\externaldocument{../05-algebrisation-partielle/algebrisation}
\externaldocument{../06-log-lissite/log}
\externaldocument{../07-uniformisation-locale/uniformisation-faible}
\externaldocument{../13-theoreme-de-finitude/finitude}
\externaldocument{../14-fonctions-dimension/fonctions-dimension}
\externaldocument{../15-Lefschetz-affine/Lefschetz-affine}
\externaldocument{../16-classes-de-Chern/gysin}
\externaldocument{../17-complexes-dualisants/dualite}
\externaldocument{../18B-dimension-cohomologique/dimensionB}
\externaldocument{../19-exemples-contre-exemples/contre-exemples}
\externaldocument{../20-rigidite/rigidite}
\externaldocument{../21-finitude-non-abelienne/finitude-non-abelienne}

\author{Luc Illusie}
\title{XVIII. Cohomological dimension: First results}

\begin{document}

\begin{center}
XVIII. Cohomological dimension: First results\\
Luc Illusie
\end{center}

\tableofcontents

\else
\setboolean{english}{true}
\setboolean{francais}{false}
\selectlanguage{english}
\chapter{Cohomological dimension: First results}
\addtocontents{toc}{par Luc Illusie}
\begin{center}
Luc Illusie
\end{center}
\fi


In this chapter we establish Gabber's bound on cohomological dimension stated in the introduction
(in the comments on the proof of the finiteness theorem).

\section{Statement of the main result and applications}

\begin{theorem}[Gabber]\label{18A-GABBER'S BOUND}
Let $X$ be a strictly local, noetherian scheme of dimension $d >0$, and let $\ell$ be a prime
number invertible on $X$. Then, for any open subset $U$ of $X$, we have
\begin{equation}\label{18A-Gabber's bound1} \mathrm{cd}_{\ell}(U) \le 2d-1.  \end{equation}
\end{theorem}
Recall that, for a scheme $S$, $\mathrm{cd}_{\ell}(S)$ ($\ell$-\textit{cohomological
dimension} of $S$) denotes the infinimum of the integers $n$ such that for all $\ell$-torsion abelian
sheaves $F$ on $S$, and all $i > n$, $H^i(S,F) = 0$.  \begin{corollary}\label{18A-fraction field}
Let $X = \SP \,A$ be as in \ref{18A-GABBER'S BOUND}, and assume $A$ is a domain, with fraction field
$K$. Then \begin{equation}\label{18A-fraction field1} \mathrm{cd}_{\ell}(K) \le 2d-1.
\end{equation} \end{corollary} Indeed, it suffices to show that if $F$ is a finitely generated
$\mathbf{F}_{\ell}$-module over $\eta = \SP\, K$, then $H^i(\eta,F) = 0$ for $i > 2d-1$. But $\eta$
is a filtering projective limit of affine open subsets $U_{\alpha}$ of $X$, $F$ is induced from a
locally constant constructible $\mathbf{F}_{\ell}$-sheaf $F_{\alpha_0}$ on $U_{\alpha_0}$, and
$H^i(\eta,F) = \varinjlim H^i(U_{\alpha},F_{\alpha})$, where $F_{\alpha} = F_{\alpha_0} |
U_{\alpha}$ for $\alpha \ge \alpha_0$ (\sga{4}{vii}{5.7}).

\begin{remark}\label{18A-fraction field
sharp} (a) The proof shows that if (\ref{18A-Gabber's bound1}) holds for one $X$, and if $X$ is
integral, then (\ref{18A-fraction field1}) holds for $X$.
 
 (b) If $X$ is an integral noetherian scheme of dimension $d$, with generic point $\SP \, K$, and
$\ell$ is a prime number invertible on $X$, then $\mathrm{cd}_{\ell}(K) \ge d$ (\sga{4}{x}{2.5}).
Gabber can prove that under the assumptions of \ref{18A-fraction field} one has
$\mathrm{cd}_{\ell}(K) = d$.
\end{remark}

\begin{corollary}\label{18A-cd(f)} Let $Y$ be a
noetherian scheme of finite dimension, $f : X \ra Y$ a morphism of finite type, and $\ell$ a prime
number invertible on $Y$. Then $$ \mathrm{cd}_{\ell}(Rf_*) < \infty, $$ i. e. there exists an
integer $N$ such that for any $\ell$-torsion abelian sheaf $F$ on $X$, $R^qf_*F = 0$ for $q > N$.
\end{corollary}

\textit{Proof of \ref{18A-cd(f)}}. We may assume $Y$ affine. Covering $X$ by finitely many open
affine subsets $U_i$ ($0 \le i \le  n)$, and using the alternate \v Cech spectral sequence $$
E^{pq}_1 = \oplus R^q(f|U_{i_0 \cdots i_p})_*(F|U_{i_0 \cdots i_p}) \Rightarrow R^{p+q}f_*F, $$
where $U_{i_0 \cdots i_p} = U_{i_0} \cap \cdots \cap U_{i_p}$, we may assume $f$ separated.
Repeating the procedure, we may assume $X$ affine. Choose an immersion $X \ra \mathbf{P}^n_Y$, and
replace $\mathbf{P}^n_Y$ by the closure of the image of $X$. We get a commutative diagram $$
\xymatrix{X \ar[d]^f \ar[r]^j & \overline{X} \ar[dl]^g \\ Y} $$ with $j$ open and $g$ projective of
relative dimension $\le n$. By the proper base change theorem we have $\mathrm{cd}_{\ell}(Rg_*) \le
2n$. By the Leray spectral sequence $R^pg_*R^qj_*F \Rightarrow R^{p+q}f_*F$ it thus suffices to
prove \ref{18A-cd(f)} for $j$, in other words, we may assume that $f$ is an open immersion. Let $d$
be the dimension of $Y$. Let $x$ be a geometric point of $X$, with image $y = f(x)$ in $Y$, and let
$U = Y_{(y)} \times_Y X$ be the corresponding open subset of the strictly local scheme $Y_{(y)}$ (of
dimension $\le d$, that we may assume to be $>0$). Then $$ R^if_*(F)_y = H^i(U,F) $$ (where we still
denote by $F$ its inverse image on $U$). The conclusion follows from \ref{18A-GABBER'S BOUND}.
\begin{remark} \label{18A-remark-d-2d-1} (a) Under the assumptions of \ref{18A-GABBER'S BOUND}, if
$X$ is quasi-excellent and $U$ is affine, then by Gabber's affine Lefschetz theorem
(\refext{XV}{15-prop-devissage-classique}) we have $\mathrm{cd}_{\ell}(U) \le d$. 

(b) Gabber can show that, under the assumptions of \ref{18A-GABBER'S BOUND}, one has
$\mathrm{cd}_{\ell}(U) \ge d$ if $U$ is not empty and does not contain the closed point and that for
each $n$ such that $d \le n \le 2d-1$, there exists a pair $(X,U)$ as in \ref{18A-GABBER'S BOUND},
with $U$ affine, such that $\mathrm{cd}_{\ell}(U) = n$ (by (a), for $n > d$, $X$ is not
quasi-excellent). These results, as well as the one mentioned in \ref{18A-fraction field sharp}, are
proved in the second part of this chapter, \emph{raffinements et compléments}, by Fabrice Orgogozo.
\end{remark}

\section{Proof of \ref{18A-GABBER'S BOUND}} \begin{lemma}\label{18A-punctured spectrum} Let $X$ be
as in \ref{18A-GABBER'S BOUND}, and let $\{x\}$ be the closed point of $X$. Then (\ref{18A-Gabber's
bound1}) holds for $U = X - \{x\}$.  \end{lemma} \textit{Proof.} It suffices to show that for any
constructible $\mathbf{F}_{\ell}$-sheaf $F$ on $U$, $H^i(U,F) = 0$ for $i \ge 2d$. Let $\widehat X$
be the completion of $X$ at $\{x\}$ and set $\widehat U := \widehat X \times_X U = \widehat X -
\{x\}$. Let $\widehat F$ be the inverse image of $F$ on $\widehat U$. By Gabber's formal base change
theorem (\cite[6.6.4]{TUBULAR@Fujiwara}), the natural map $$ H^i(U,F) \ra H^i(\widehat U, \widehat
F) $$ is an isomorphism for all $i$. Therefore we may assume $X$ complete, and in particular,
excellent. Let $(f_1,\cdots,f_d)$ be a system of parameters of $X$, and let $U_i = X_{f_i}$, so that
$U = \cup_{1 \le i \le d}U_i$. Consider the (alternate) \v Cech spectral sequence $$ E^{pq}_1 =
\oplus H^q(U_{i_0 \cdots i_p},F) \Rightarrow H^{p+q}(U,F), $$ with $U_{i_0 \cdots i_p} = U_{i_0}
\cap \cdots \cap U_{i_p}$ as above.  By definition, $E^{pq}_1 = 0$ for $p \ge d$. On the other hand,
as $X$ is excellent, by Gabber's affine Lefschetz theorem
(\refext{XV}{15-prop-devissage-classique})), we have $E^{pq}_1 = 0$ for $q \ge d+1$. Therefore
$E^{pq}_1 = 0$ for $p+q \ge 2d$, hence $H^i(U,F) = 0$ for $i \ge 2d$.
\begin{lemma}\label{18A-punctual bound} Let $X$ be a noetherian scheme of finite dimension, $Y$ a
closed subset, $\ell$ a prime number invertible on $X$. Then, for any $\ell$-torsion sheaf $F$ on
$X$, $$ H^i_Y(X,F) = 0 $$ for $$ i > \mathrm{sup}_{x \in Y} (\mathrm{cd}_{\ell}(k(x)) + 2
\mathrm{dim} \OO_{X,x}).  $$ In particular, $$ \mathrm{cd}_{\ell}(X) \le \mathrm{sup}_{x \in
X} (\mathrm{cd}_{\ell}(k(x)) + 2 \mathrm{dim} \OO_{X,x}).  $$ \end{lemma} \textit{Proof.}
For $p \ge 0$, let $\Phi^p$ be the set of closed subsets of $Y$ of codimension $\ge p$ in $X$. We
have $\Phi^p = \emptyset$ for $p > \mathrm{codim}(Y)$.  Consider the (biregular) coniveau spectral
sequence of the filtration $(\Phi^p)$ cf. \cite[10.1]{BrauerIII@Grothendieck}),
\begin{equation}\label{18A-coniveau} E^{pq}_1 = H^{p+q}_{\Phi^p/\Phi^{p+1}}(X,F) \Rightarrow
H^{p+q}_Y(X,F).  \end{equation} We have $$ E^{pq}_1 = \oplus_{x \in Y^{(p)}}
H^{p+q}_{\{x\}}(X_x,F|X_x), $$ where $Y^{(p)}$ denotes the set of points of $Y$ of codimension $p$
in $X$, and $X_x = \SP \, \OO_{X,x}$.  For $x \in Y^{(p)}$ (i. e.
$\mathrm{dim}\OO_{X,x} = p$), let $\overline{x}$ be a geometric point above $x$. Consider
the diagram $$ \xymatrix{\{\overline{x}\} \ar[d] \ar[r]^{i_{\overline{x}}} & X_{(\overline{x})}
\ar[d] & \overline{U} \ar[l]_{\overline{j}} \ar[d] \\ \{x\} \ar[r]^{i_x} & X_x & U \ar[l]_{j}}, $$
where $U = X - \{x\}$, $\overline{U} = X_{(\overline{x})}- \{\overline{x}\}$. We have $$
R\Gamma_{\{x\}}(X_x,F|X_x) = R\Gamma(\{x\},Ri_x^!(F|X_x)).  $$ The stalk of $Ri_x^!(F|X_x)$ at
$\overline{x}$ is $$ Ri_x^!(F|X_x)_{\overline{x}} = Ri^!_{\overline{x}}(F| X_{(\overline{x})}), $$
as $(Rj_*(F|U))_{\overline{x}} = R{\overline{j}}_*(F|\overline{U})_{\overline{x}}$.  We thus have a
spectral sequence \begin{equation}\label{18A-Galois} E^{rs}_2 = H^r(k(x),R^si_{\overline{x}}^!(F|
X_{(\overline{x})})) \Rightarrow H^{r+s}_{\{x\}}(X_x,F|X_x) , \end{equation} It suffices to show
that, in the initial term of (\ref{18A-coniveau}), \begin{equation}\label{18A-E1}
H^{p+q}_{\{x\}}(X_x,F|X_x) = 0 \end{equation} for $p+q > \mathrm{cd}_{\ell}(k(x)) + 2p$. If $p = 0$,
then $Ri^!_{\overline{x}}(F| X_{(\overline{x})}) = F_{\overline{x}}$, and, in (\ref{18A-Galois}),
$E^{rs}_2 = 0$ for $s >0$, $E^{r0}_2 = 0$ for $r > \mathrm{cd}_{\ell}(k(x))$, so (\ref{18A-E1}) is
true in this case. Assume $p >0$. We have \begin{equation}\label{18A-local cohomology}
R^si_{\overline{x}}^!F = H^{s-1}(\overline{U},F|\overline{U}) \end{equation} for $s \ge 2$, where,
as above, $\overline{U} = X_{(\overline{x})} - \{\overline{x}\}$. By \ref{18A-punctured spectrum},
$H^{s-1}(\overline{U},F| \overline{U}) = 0$ for $s-1 \ge 2p$, hence, by (\ref{18A-local
cohomology}), $R^si_{\overline{x}}^!(F|X_{(\overline{x})}) = 0$ and $E^{rs}_2 = 0$ for $s \ge 2p
+1$. If $r+s \ge \mathrm{cd}_{\ell}(k(x)) + 2p + 1$ and $s \le 2p$, then $r >
\mathrm{cd}_{\ell}(k(x))$, hence $E^{rs}_2 = 0$ as well. Therefore, by (\ref{18A-Galois}),
(\ref{18A-E1}) holds, which finishes the proof. 

\medskip \textit{Proof of \ref{18A-GABBER'S BOUND}}. We prove \ref{18A-GABBER'S BOUND} by induction
on $d$. For $n \ge 0$ consider the assertion

\medskip $G_n$ : \textit{For every strictly local, noetherian scheme $X$ of dimension $n$, all open
subsets $U$ of $X$ and any prime number $\ell$ invertible on $X$, we have $\mathrm{cd}_{\ell}(U) \le
\mathrm{sup}(0,2n-1)$}. 

\medskip Let $d >0$. Assume $G_n$ holds for $n < d$, and let us prove $G_d$. Let $X$ be as in
\ref{18A-GABBER'S BOUND}. If $(X_i)_{1\le i \le r}$ are the reduced irreducible components of $X$
and $U_i = U \times_X X_i$, we have $\mathrm{cd}_{\ell}(U) \le \mathrm{sup}(
\mathrm{cd}_{\ell}(U_i))$, hence we may assume $X$ integral. Let $x$ be the closed point of $X$, and
$U = X - \{x\}$ the punctured spectrum. Let $j : V \ra U$ be a nonempty open subset of $U$, and $F$
be a constructible $\mathbf{F}_{\ell}$-sheaf on $V$. As $F = j^*j_!F$, by \ref{18A-punctured
spectrum} it suffices to show that, for any constructible $\mathbf{F}_{\ell}$-sheaf $L$ on $U$, the
restriction map $$ H^i(U,L) \ra H^i(V,j^*L) \leqno (*) $$ is an isomorphism for $i \ge 2d$. Let $Y =
U - V$. Consider the exact sequence $$ H^i_Y(U,L) \ra H^i(U,L) \ra H^i(V,j^*L) \ra H^{i+1}_Y(U,L).
$$ By \ref{18A-punctual bound}, we have $H^i_Y(U,L) = 0$ for $i > \mathrm{sup}_{y \in Y}
(\mathrm{cd}_{\ell}(k(y)) + 2 \mathrm{dim} \OO_{X,y})$. For $y \in Y$, denote by $Z$ the
closed, integral subscheme of $X$ defined by the closure of $\{y\}$ \textit{in} $X$. As $X$ is
integral and $V$ nonempty, $Z$ is a strictly local scheme of dimension $n < d$, with generic point
$y$.    By \ref{18A-fraction field sharp} (a) and $G_n$ (inductive assumption), we have
$\mathrm{cd}_{\ell}(k(y)) \le 2n -1$. We have $2n - 1 + 2\mathrm{dim}\OO_{X,y} \le 2d-1$.
Hence, for $i \ge 2d$, $H^i_Y(U,L) = H^{i+1}_Y(U,L) = 0$, and (*) is an isomorphism, which finishes
the proof.  \begin{remark} Gabber has an alternate proof of \ref{18A-GABBER'S BOUND}, based on the
theory of Zariski-Riemann spaces. By \ref{18A-punctual bound}, it suffices to show \ref{18A-fraction
field}. Here is a sketch, pasted from an e-mail of Gabber to Illusie of 2007, Aug. 15 :

\medskip "For $Y \ra X$ proper birational with special fiber $Y_0$, consider $i: Y_0 \ra Y$ and $j
: \eta \ra Y$, $\eta$ the generic point. We have by proper base change a spectral sequence $$
H^p(Y_0,i^*R^q j_*F) \rightarrow H^{p+q}(\eta,F) $$ for $F$ an $\ell$-torsion Galois module.  We
take the direct limit and get a spectral sequence involving cohomologies on the étale topos of
$ZRS_0$ defined as the limit of étale topoi of $Y_0$ or viewing $ZRS_0$ as a locally ringed topos
and applying a universal construction in the book of M.Hakim. The limit of the $R^q j_* F$ is $R^q
(\eta \ra ZRS)_* F$.  By a classical result of Abhyankar, also proved in Appendix 2 of the book of
Zariski-Samuel Vol. II, if $R$ is a noetherian local domain of dimension $d$ and $V$ a valuation
ring of $\mathrm{Frac}(R)$ dominating $R$, the sum of the rational rank and the residue
transcendence degree is at most $d$.  For a strictly henselian valuation ring  $V$ with residue
characteristic exponent $p$ and value group $\Gamma$, the absolute Galois group of
$\mathrm{Frac}(V)$ is an extension of the tame part (product for $\ell$ prime not equal to $p$ of
$\mathrm{Hom}(\Gamma, \mathbf{Z}_{\ell}(1)))$ by a $p$-group, so the $\ell$-cohomological dimension
is the dimension of $\Gamma$ tensored with the prime field $\mathbf{F}_{\ell}$, which is at most the
dimension of $\Gamma$ tensored with the rationals.  If $A$ is an $\ell$-torsion sheaf on the étale
topos of $ZRS_0$, let $\delta(A)$ be the sup of transcendence degrees of points where the stalk is
non-zero. I claim that $H^n(ZRS_0,A)$ vanishes for $n> 2 \delta(A)$. One reduces it to the finite
type case (passage to the limit \sga{4}{vi}{8.7.4})  using that the $\delta$ of the direct image of
$A$ to $Y_0$ is at most $\delta(A)$.  In $Y_0$ the transcendence degrees over the closed point of
$X$ are at most $d-1$ by the dimension inequality. Summing up, for the limit spectral sequence the
$q$-th  direct image sheaf  restricted to the special fiber has $\delta$ at most
$\mathrm{min}(d-1,d-q)$, giving vanishing for certain $E_2^{p,q}$ and the result." \end{remark}

\ifx\danslelivre\undefined
\bibliography{../fichiers-configuration-latex/bibliographie}
\bibliographystyle{../fichiers-configuration-latex/smfalpha}
\end{document}
\fi
